\input header
\input ras.tex

\footline={\hss\tenrm 17-\folio\hss}\pageno=1

\centerline{\bf POLYNOMIAL MAPS}\bigskip

\centerline{\bf \S17. The Filled Julia Set.}\bigskip

Let \[f:\hat\C\to\hat\C\] be a polynomial map of degree \[d\ge 2\], say
$$	f(z)=a_dz^d+\cdots+a_1z+a_0	$$
with \[a_d\ne 0\]. Then \[f\] has a superattracting fixed point at infinity.
In particular, there exists a constant \[k=k_f\] so that any point \[z\]
with \[|z|>k_f\] belongs to the basin of attraction of infinity.

{\bf Definition.} The complement of the attractive basin of infinity, that
is the set of all points \[z\in\C\] with bounded forward orbit under \[f\],
is called the {\bit filled Julia set} \[K=K(f)\].

{\QP{\bf 17.1. Lemma.} {\it This filled Julia set \[K\] is a compact set
consisting of Julia set \[J\]
together with all of the bounded components of the complement \[\C\ssm J\].
These bounded components $($if any$)$ are all simply connected.
The Julia set \[J\] is equal to the topological boundary \[\partial K\].
}\medskip}

{\bf Proof.} (Compare Problem 4-1.)
If \[z\] is a boundary point of \[K\], then \[z\] itself has
bounded orbit, but points arbitrarily close have orbits which converge to
the point at infinity. Thus the family of iterates of \[f\] cannot be normal
in any neighborhood, hence \[z\in J\]. Conversely, if \[z\in J\], then
\[z\] has bounded orbit, but it follows from 4.6 that points arbitrarily
close to \[f\] have unbounded orbit, which therefore converges to infinity.
Thus \[J=\partial K\subset K\], and it follows that the unbounded component
of the complement \[\C\ssm J\] consists of points attracted to infinity. But if
\[z\] belongs to a bounded component of \[\C\ssm J\], then it
follows from the Maximum Modulus Principle that every point in the orbit
of \[z\] belongs to the closed disk of radius \[k_f\], hence \[z\in K\].
Finally, any bounded component \[U\] of \[\C\ssm J\] must be simply connected,
since for any Jordan
curve \[\Gamma\subset U\] the entire bounded component of \[\C\ssm \Gamma\] must
also be contained in \[U\], again by the Maximum Modulus Principle.
\QED

It follows from the B\"ottcher
Theorem \S6.7 that exists a new coordinate \[w=\phi(z)\]
which is defined throughout some neighborhood \[|z|>c_f\] of infinity, with
\[w\to\infty\] as \[z\to\infty\], and which satisfies
$$	\phi\circ f\circ\phi^{-1}(w)\;=\; w^d	\,. $$
Further, \[\phi\] is unique up to multiplication with a \[(d-1)$-st root
of unity.\smallskip

The function \[G(z)=\log|\phi(z)|\] is called the canonical {\bit potential
function} or {\bit Green's function}$\,$ for \[K(f)\]. Note the identity
$$	G(z)\;=\;G(f(z))/d\,.	\eqno (*)	$$
Although we have so far defined \[G\] only in a neighborhood of infinity,
there is one and only one extension to all of \[\C\] which is continuous
and satisfies (*). In fact we define \[G(z)=0\] for \[z\in K(f)\], and
\[G(z)=G(f^{\circ n}(z))/d^n\] otherwise, where \[n\] is large enough so
that \[|f^{\circ n}(z)|>c_f\]. Alternatively, we can simply set
$$	G(z)\;=\;\lim_{n\to\infty}\, \log^+|f^{\circ n}(z)|/d^n\,, $$
where \[\log^+(x)=\log(x)\] for \[x\ge 1\] and \[\log^+(x)=0\] for
\[0\le x\le 1\].
Note that \[G\] is a smooth real analytic function outside of the Julia set
\[J\], but is only continuous at points of \[J\].
The verification of these facts will be left to the reader.

The lines
\[|\phi(z)|={\rm constant}>0\] are called {\bit equipotential curves}.
Their orthogonal trajectories, that is the images under \[\phi^{-1}\]
of the radial lines extending out towards infinity from the unit disk
are called {\bit external rays} for \[f\]. Both of these families of
curves are clearly
smooth except at critical points of \[G\]. {\it Note that a point \[z\in
\C\ssm K\] is a critical point of \[G\] if and only if it is a {\bit pre-critical
point} of \[f\], that is a critical point of some iterate \[f\circ\cdots
\circ f\].} Again proofs are easily supplied.

For each \[r>1\] let \[V(r)\] be the compact region with boundary consisting
of all \[z\] with \[G(z)\le\log(r)\]. Evidently the boundary \[\partial
V(r)\] is just the equipotential curve \[\{z:G(z)=\log r\}\].

{\QP{\bf 17.2. Lemma.} \it Suppose that there are no critical points of \[G\]
$($or pre-critical points of \[f\,)\] on the
boundary \[\partial V(r)\]. If \[m\ge 0\] is the algebraic number of critical
points of \[G\] in the complement \[\C\ssm V(r)\], then \[V(r)\] is a disjoint
union of \[m+1\] closed topological disks, each of which intersects the
Julia set \[J(f)\].\medskip}

{\bf Proof.} If \[\partial V(r)\] contains no critical points (or equivalently
if \[\log r\] is a regular value of the map \[G\]), then evidently \[V(r)\]
is a smooth compact manifold with boundary. Each component must be a
closed topological disk. For otherwise, the complement \[\C\ssm V(r)\] would
have a bounded component, which is impossible since the function \[G\]
would have to take its maximum on the boundary of such a component. Similarly,
the minimum value of \[G\] on each component of \[V(r)\] must occur at
a point of \[K(f)\]. Hence each such component must intersect the Julia set
\[\partial K(f)\].

Given \[r>1\] as above, let us choose \[n\] so large that the region
\[V(r^{d^n})\] is a single topological disk. It will be convenient to
set \[s=r^{d^n}\]. Then the composition \[f^{\circ n}
\] maps \[V(r)\] onto \[V(s)\] by a \[d^n$-fold branched covering map.
According to the Riemann-Hurwitz formula, \S5.1, the algebraic number
of branch points of \[f^{\circ n}\] in \[V(r)\] is equal to \[d^n\chi(V(
s))-\chi(V(r))=d^n-\chi(V(r))\],
where in our case \[\chi\] is just the number of
connected components. If \[m\] is the number of branch points of \[f^{\circ n}
\] outside of \[V(r)\], then, since the total number of branch points is
\[d^n-1\], we can write this formula as
$$	d^n-1-m\;=\; d^n-\chi(V(r))\,,	$$
or in other words \[\chi(V(r))=m+1\]. This completes the proof,
since critical points of \[G\] are the same as critical points of
\[f^{\circ n}\] for large \[n\].\QED

Following Brown, a compact set in Euclidean space is said to be$\,$ {\bit
cellular} $\,$ if it is equal to the intersection of a nested sequence
of closed embedded topological disks, each of which contains the next
in its interior. The following result is essentially due to Fatou and Julia.

{\QP{\bf 17.3. Theorem.} {\it For a polynomial map \[f\] of degree \[d\ge 2\],
there are just two mutually exclusive possibilities: If the
filled Julia set \[K\] contains all of the finite critical points of \[f\],
then \[K\] and \[J=\partial K\] are connected, and in fact \[K\] is a cellular
set. Furthermore, the B\"ottcher map near infinity extends to a conformal isomorphism
$$	\phi\;:\;\C\ssm K\;\;\buildrel\approx\over\longrightarrow\;\; \C\ssm\bar D\,.	$$
On the other hand, if \[f\]
has at least one critical point in \[\C\ssm K\], then both \[J\] and \[K\]
have uncountably many connected components.}\medskip}

{\bf Proof of 17.3.} If \[K\] contains all of the finite critical points,
then it follows from 17.2 that \[K\] is the intersection of a strictly nested
family of embedded disks \[V(r)\]. Thus \[K\] is cellular, and hence connected.
Furthermore, the proof shows that each \[\C\ssm V(r)\] is a \[d^n$-fold unramified
covering of some \[\C\ssm V(s)\], which is isomorphic to \[\C\ssm\bar D_s\] under
the B\"ottcher map. We can then lift to a B\"ottcher isomorphism
\[\C\ssm V(r)\buildrel\approx\over\to\C\ssm\bar D_r\] which is compatible near
infinity and satisfies
$$	\phi(f(z))=\phi(z)^d\,.	$$
Passing to the limit as \[r\to 1\], we see that \[\phi\] can be defined
throughout \[\C\ssm K\]. Finally, note that the Julia set
\[J\] can be expressed as the intersection of the closures of the bounded
connected annuli \[\phi^{-1}(\{w:1<|w|<r\})\]; hence \[J\] is also connected.

Conversely, if \[f\] has at least one critical point in \[\C\ssm K\], then the
argument above shows that some \[V(r)\] has two or more connected components,
and that each of these components
intersects the subsets \[J\subset K\subset V(r)\]. Thus \[J\] has uncountably
many components by 11.3. The analogous proof for \[K\] will be left to the
reader.\QED

For information on the structure of \[K\] when it is not connected,
see Branner \& Hubbard, Blanchard.
As an immediate consequence of Carath\'eodory's Theorem 16.6
we have the following.
(Note that a compact set which has infinitely many components clearly cannot
be locally connected.)

{\QP{\bf 17.4. Corollary.} {\it With \[f\] as above, the following three conditions
are equivalent:

$(a)$ The set \[J(f)\] is locally connected.

$(b)$ The set \[K(f)\] is locally connected.

$(c)$ K(f) is connected, and the inverse Riemann mapping \[\psi:\hat\C\ssm\bar
D\buildrel\approx\over\to\hat\C\ssm K\] extends continuously over the boundary,
yielding a continuous map from the unit circle onto \[J\] satisfying
the identity \[f(\psi(w))=\psi(w^d)\].}\medskip}

\noindent Following Douady and Hubbard, we have the following.

{\QP{\bf 17.5. Theorem.} {\it If the polynomial map \[f\] is hyperbolic or
subhyperbolic,
with \[K(f)\] connected, then \[K(f)\] is locally connected.}\medskip}

For non locally connected examples, see 18.6.\smallskip

{\bf Proof of 17.5.} We may assume that there is a smooth Riemannian metric
or orbifold metric
whose restriction to some small neighborhood \[U\] of \[J\] is strictly
increasing under \[f\],
say by a factor of at least \[k>1\]. Choose \[r>1\] small enough so that
the neighborhood \[V(r)\] of \S17.2 is contained in \[K\cup U\], and let
\[s=\root d\of r\]. Consider the (half-closed) annulus
$$	A_0\;=\; V(r)\ssm V(s)\,,	$$
which is isomorphic under \[\phi\] to the ``round'' annulus \[\bar D_r\ssm\bar
D_s\]. Then \[V(r)\ssm K\] can be expressed as the countable
union \[A_0\cup A_1\cup
\cdots\], where \[A_n\] is the annulus \[f^{-n}(A_0)\]. Similarly,
\[\bar D_r\ssm\bar D_1\] is the union of a corresponding family of round
annuli.

Each radial line segment in \[\bar D_r\ssm\bar D_s\] corresponds
to a smooth curve in \[A_0\], which is a segment of some external ray.
Let \[\ell_0\] be the
maximum of the lengths of these curve segments in \[A_0\], and let \[\ell_n\] be the
maximum length of the corresponding curve segments in \[A_n\]. Since \[f^{\circ n}\]
maps \[A_n\] onto \[A_0\], preserving this foliation by external rays,
and stretching all distances by at least \[k^n\], it follows that
\[\ell_n\le\ell_0/k^n\]. Therefore, the total length of any external ray
intersected with \[V(r)\] is uniformly bounded by the quantity\break \[\;\ell_0
(1+k^{-1}+k^{-2}+\cdots)\;=\;\ell_0
/(k-1)\;<\;\infty\,\]. It then follows easily that the maps
$$	e^{i\theta}\;\mapsto\;\psi(re^{i\theta})	$$
from the unit circle onto \[\partial V(r)\]
converge uniformly as \[r\to 1\] to the required limit mapping, which
carries the unit circle continuously onto the Julia set.\QED
\vfil\eject % \bigskip\bigskip

\footline={\hss\tenrm 18-\folio\hss}\pageno=1


\centerline{\bf \S18. External Rays and Periodic Points.}\bigskip

To fix our ideas and simplify the discussion, let us assume
that the filled Julia set \[K\] is connected. (For the general case, see
[DH2].) By 17.3 the complement \[\C\ssm K\]
is isomorphic to the complement \[\C\ssm\bar D\] under a conformal isomorphism
$$	\phi:\C\ssm K\;\buildrel\approx\over\longrightarrow\;\C\ssm\bar D	$$
which is compatible with the dynamics:
$$	\phi(f(z))\;=\;\phi(z)^d\,.	$$
Here \[\phi\] is unique up to multiplication by a \[(d-1)$-st root of unity.
In practice, it is customary to make a linear change of coordinates so that
the polynomial \[f\] is {\bit monic}, ie., has leading coefficient equal to
\[+1\]. There is then a preferred choice of B\"ottcher coordinate \[w=
\phi(z)\], determined by the requirement that \[\;w/z\to 1\] as \[|z|\to
\infty\].\smallskip

{\bf Definition.} Let \[\psi:\C\ssm\bar D\buildrel\approx\over\to\C\ssm K\] be the inverse map.
The image under \[\psi\] of a radial line
$$	\{re^{2\pi it}:r>1\}	$$
in \[\C\ssm\bar D\] is called the {\bit external ray} \[\;R_t\]
at {\bit angle} \[t\] in \[\C\ssm K\]. Note that our {\it angles} are elements of
\[\R/\Z\], measured in fractions of a full turn and not in radians. Evidently
$$	f(R_t)\;=\;R_{dt}	\,.	$$
That is: {\it \[f\] maps the external ray at angle \[t\] to the external ray
at angle \[dt\].} In particular, if \[t\] is a fraction of the form \[m/
(d-1)\] then \[f(R_t)=R_t\].\smallskip

By definition, the external ray \[R_t\] {\bit lands} at some point
\[z_t\], which necessarily belongs to the Julia set \[J\], if the points
\[\psi(re^{2\pi it})\in \C\] tend to a well defined limit \[z_t\] as \[r\to
1\]. If \[J\], or equivalently \[K\], is locally connected, then according to
17.4 {\it every} external ray lands, at a point \[z_t\] which depends
continuously on \[t\]. Douady and Hubbard call the resulting parametrization
\[t\mapsto z_t\] the {\bit Carath\'eodory loop} on the Julia set.
However, in this section we will usually not assume that \[J\] is locally
connected.\smallskip

{\bf Remark.} A priori, it is possible that every external ray may land
even if \[K\]
is not locally connected. An example of a compact set (but not a filled Julia
set) with this property is shown in Figure 16. Evidently, in such a case,
the correspondence \[t\mapsto z_t\] cannot be continuous.\smallskip
\midinsert
\insertRaster comb.ps pixels 512 by 513 scaled 225
\bigskip
\centerline{Figure 16. The double comb \[\;[0,2]\times\{0\,,\,\pm 2^{-n}\}\;
\cup\;0\times[-1,1]\]}
\medskip\vfil
\endinsert
%\eject

{\bf Definition.} An external ray \[R_t\] will be called {\bit rational}
if its angle \[t\in\R/\Z\] is rational, and {\bit periodic} if
\[t\] is periodic under multiplication by the degree \[d\],
so that \[\;d^nt\equiv t \;(\mod 1)\] for some \[n\ge 1\].\smallskip

Note that \[R_t\] is eventually periodic
under  multiplication by \[d\] if and only if \[t\] is rational, and
is periodic if and only if the number \[t\] is
rational with denominator relatively prime to \[d\].
(If \[t\] is rational with denominator \[m\], then
the successive images of
\[R_t\] under \[f\] have angles \[dt\,,\,d^2t\,,\,d^3t\,,\,\ldots\;
{\rm (mod\;1
)}\;\;\] with denominators dividing \[m\]. Since there are only finitely
many such fractions modulo 1, this sequence must eventually repeat.\break
% Using ???, it follows that the corresponding sequence
% $$	z_t\buildrel f\over\mapsto z_{dt}\buildrel f\over\mapsto z_{d^2t}
%	\buildrel f\over\mapsto\cdots	$$
%is also eventually periodic. 
In the special case where \[m\] is relatively
prime to \[d\], the fractions with denominator \[m\] are permuted under
multiplication by \[d\] modulo 1, so the landing point \[z_t\] is
actually periodic.)\smallskip

%(Conversely, if \[z_t\] is periodic, then the denominator \[m\]
%must be relatively prime to \[d\]. For otherwise we could find
%two distinct fractions of the form \[d^nt\] modulo 1 which have the same
%image under multiplication by \[d\] modulo 1. Hence we could find
%two distinct external rays landing at
%points on the orbit of \[z_t\], with a common image under \[f\]. But
%this is impossible: Since \[z_t\] is in the Julia set, it cannot belong to
%a superattractive orbit. Therefore, \[f\] is a local diffeomorphism near the
%orbit of \[z_t\].)\smallskip

The following result is due to Sullivan, Douady and Hubbard.
We do not assume local connectivity.

{\QP{\bf 18.1. Theorem.} \it If \[K(f)\] is connected, then every periodic
external ray lands at a periodic point which is either repelling or
parabolic.\medskip}

\noindent The converse result, due to Douady and Yoccoz, is more difficult:

{\QP{\bf 18.2. Theorem.} \it Still assuming that \[K(f)\] is connected,
every repelling or parabolic periodic point is the landing point of at
least one external ray, which is necessarily periodic.\medskip}

In fact we will only prove the parabolic case. For the repelling case, which
is quite similar, the reader is referred to Petersen.
Here is a complementary result.

{\QP{\bf 18.3. Lemma.} \it If a periodic ray lands at \[z_t\], then only
finitely many external rays, all periodic of the same period, can land
at \[z_t\].\medskip}

\noindent The following is an immediate consequence of 18.1.

\midinsert
\vfil
\insertRaster cubicj-rays.ps pixels 480 by 960 scaled 350
\bigskip
\centerline{Figure 17. Julia set for \[z\mapsto z^3-iz^2+z\]}
\centerline{with some external rays and equipotentials indicated.}
\vfil
\bigskip
\endinsert

{\QP{\bf 18.4. Corollary.} \it If \[t\] is rational but not periodic,
then \[R_t\] lands at a point \[z_t\] which is eventually periodic
but not periodic.\medskip}

As an example, Figure 17 shows the Julia set for the cubic map \[f(z)=z^3-
iz^2+z\].
Here the rays \[R_0\] and \[R_{1/2}\] are mapped into themselves by \[f\].
Both must land at the parabolic fixed point \[z=0\], since the only other
fixed point, at \[z=i\], is superattracting and hence does not belong
to the Julia set. On the other hand,
the \[1/6\,,\,1/3\,,\, 2/3\] and \[5/6\] rays have denominator divisible
by 3. These four rays land at the two disjoint pre-images of zero.\smallskip

The following result of Sullivan and Douady is closely related.
(Sullivan 1983. For an independent proof, see Lyubich, p. 85.)

{\QP\bf 18.5. Theorem. {\it If the Julia set \[J\] of a polynomial map \[f\]
is locally connected, then every periodic point in \[J\] is either
repelling or parabolic.}\medskip}

Recall from \S8
that a {\bit Cremer point} can be characterized as a periodic point which
belongs to the Julia set but is neither repelling nor parabolic. Thus
the following is a completely equivalent formulation.

{\QP\bf 18.6. Corollary. {\it If \[f\] is a polynomial map with a Cremer
point, then the Julia set \[J(f)\] is not locally connected.}\medskip}

\noindent The proofs begin as follows.

{\QP{\bf 18.7. Lemma.} {\it The ray \[R_t\] lands at a single point \[z_t\]
of the Julia set if and only if \[R_{dt}\] lands at a single point
\[z_{dt}\]. Furthermore, the image \[f(z_t)\] is necessarily equal to
\[z_{dt}\].}\medskip}
\eject

{\bf Proof.} The set \[L(t)\] consisting of all limit points of
\[\psi(re^{2\pi it})\] as \[r\] tends to 1 from above
is certainly compact and connected, with \[f(L(t))=L(dt)
\].  It follows easily that \[L(t)\] is a single point if and only if
\[L(dt)\] is a single point.\QED

Thus if the ray \[R_t\] is periodic, and if \[R_t\] lands at a point \[z_t\],
then it follows that \[z_t\] is a periodic point of \[f\].\smallskip

{\bf Proof of 18.3.}
First consider the special case \[t=0\]. Then \[f\] maps the ray \[R_0\] to
itself, hence \[f(z_0)=z_0\]. If another ray \[R_s\] also lands at \[z_0\],
then we will prove that \[s\] is a rational number of the form \[j/(d-1)\].
This will show that there are
at most \[d-1\] external rays landing at \[z_0\], and that they
are all periodic (and in fact fixed) under multiplication by \[d\], since
the numbers \[s=j/(d-1)\] are exactly the solutions to the congruence
\[s\equiv ds\;({\rm mod}\; 1)\] where \[d\] is the degree.

If a ray \[R_s\] which is not of
this form lands at \[z_0\], then we have \[s\not\equiv ds\;({\rm mod}\; 1)\].
Define a sequence of angles \[0\le s_n<1\] by the condition that \[s_0\equiv
s\] and \[s_{n+1}\equiv ds_n\] modulo 1. By hypothesis, the numbers \[0\,,
\,s_0\] and \[s_1\] are distinct.
 Suppose, to fix our ideas, that
\[0<s_0<s_1<1\]. Since \[f\] is a local diffeomorphism near \[z_0\], it must
preserve the cyclic order of the rays landing at \[z_0\]. Thus the images of
\[0\,,\,s_0\] and \[s_1\] must satisfy \[0<s_1<s_2<1\]. Continuing
inductively, it follows that \[0<s_0<s_1<s_2<\cdots<1\]. Thus the \[s_n\]
must converge to some angle \[s_\infty\], which is necessarily a fixed point
of the map \[s\mapsto ds\;({\rm mod}\; 1)\]. But this is impossible, since
this map has only strictly repelling fixed points.

More generally, if \[R_t\] is any ray which is mapped into itself by \[f\],
so that \[t=j/(d-1)\], then we can translate all angles by \[-t\] and apply
the argument above.
The proof in the general case now follows easily. We can
replace \[f\] by some iterate \[g=f^{\circ n}\], where \[d^ns\equiv s\;({\rm
mod}\; 1)\] so that \[g(R_s)=R_s\]. The argument then proceeds as above.
This proves 18.3.\QED

The proof of 18.1 begins as follows. We will make
use of the \Poincare metric on \[\C\ssm K\]. {\it Note first that the map \[f\] is
a local isometry for this metric.} In fact the
universal covering of \[\C\ssm\bar D\] is isomorphic to the right half-plane
\[\{W=U+iV:U>0\}\] under the exponential map, and the map \[f\] on \[\C\ssm K\]
corresponds to the \[d$-th power map on \[\C\ssm\bar D\], which corresponds
to the automorphism \[W\mapsto dW\] on this right half-plane.\smallskip

We suppose that \[t\] has denominator prime to \[d\], so that some
iterate \[g=f^{\circ n}\] maps \[R_t\] to itself.
This ray \[R_t\] can be expressed as a union
$$	R_t\;=\;\cdots\cup S(-1)\cup S(0)\cup S(1)\cup\cdots	$$
of finite segments, where \[g\] maps
each \[S(j)\] isomorphically onto \[S(j+1)\]. In fact let \[S(j)\] correspond
under \[\phi\] to the set of \[re^{2\pi i t}\] with \[d^{jn}\le\log r\le
d^{(j+1)n}\]. Thus the \[S(j)\] all have the same \Poincare length. But as
\[j\to-\infty\] the segments \[S(j)\] converge towards the boundary \[J\]
of \[\C\ssm K\]. Hence, according to \S2.4, the Euclidean length of \[S(j)\]
must converge to zero. This means that any limit point of the sequence
of sets \[S(j)\], as \[j\to-\infty\], must be a fixed point of \[g\].
But, as noted in the proof of 18.7,
the set of all limit points must be connected.
Since \[g\] has only finitely many fixed points, this proves that the ray
\[R_t\] lands at a single fixed point \[z_t\]
of the map \[g=f^{\circ n}\].\smallskip

Now let us apply the Snail Lemma of \S13.
After translating coordinates so that
our fixed landing point lies at the origin, the map \[g\] and the path
\[x\mapsto \psi(\exp(d^{nx}+2\pi it))\] will satisfies the hypothesis of 13.3.
Hence it follows that this fixed point \[z_t\] of \[g\] is either repelling,
or parabolic with multiplier \[g'(z_t)=1\]. In terms of the original map
\[f\], it follows that \[z_t\] is a periodic point, and that its periodic
orbit is either repelling or has a root of unity as multiplier. This proves
18.1.\QED

{\bf Proof of 18.4.} Consider a rational angle \[t\] with denominator which
is not relatively prime to \[d\]. In this case the multiple \[dt\] will
have smaller denominator,
but the sequence of angles \[d^nt\] modulo 1 must certainly be eventually
periodic. Thus some forward image \[f^{\circ n}R_t\] must have a
well defined landing point. But this implies that \[R_t\] itself has
a well defined landing point by 18.7.\QED

The proof of 18.5 and 18.6 will be based on 18.1, together with the
following. Let \[n\ge 2\] be an integer.

{\QP\bf 18.8. Lemma. {\it Let \[L\subset \R/\Z\] be a compact set which is
mapped homeomorphically onto itself by the map \[t\mapsto nt (\mod 1)\].
Then \[L\] is finite.}\medskip}

{\bf Proof.} In fact we will prove the following more general
statement. Let \[X\] be a compact metric space with distance function
\[\rho(x,y)\], and let \[h:X\to X\] be a homeomorphism which is expanding
in the following sense: There should exist numbers \[\epsilon>0\] and \[k>1\]
so that
$$	\rho(h(x),h(y))\;\ge\;k\rho(x,y)	$$
whenever \[\rho(x,y)<\epsilon\]. Then we will show that \[X\] is finite.
Evidently this hypothesis is satisfied with \[X=L\subset\R/\Z\], so
this argument will prove 18.8.

Since \[h^{-1}\] is uniformly continuous, we can choose \[\delta>0\] so that
\[\rho(x,y)<\epsilon\] whenever \[\rho(h(x),h(y))<\delta\]. But this
implies that \[\rho(x,y)<\delta/k\]. Since \[X\] is compact, we can choose
some finite number, say \[n\], of balls of radius \[\delta\] which cover
\[X\]. Applying \[h^{-p}\], we obtain \[n\] balls of radius \[\delta/k^p\]
which cover \[X\]. Since \[p\] can be arbitrarily large, this proves that
\[X\] can have at most \[n\] distinct points.\QED

{\bf Proof of 18.5.}
Let \[z_0\] be any periodic point. After
replacing the given polynomial map by some iterate, we may assume that
\[z_0\] is fixed by \[f\]. Furthermore, after a linear
change of coordinates, we may assume that \[f\] is monic.
Since \[J(f)\] is locally connected, according to 17.4 the inverse Riemann map
$$	\psi\,:\,\hat\C\ssm\bar D\;\buildrel\approx\over\longrightarrow
	\; \hat\C\ssm K(f)	$$
extends continuously over the boundary, yielding a map from the unit
circle onto \[J(f)\]
which semi-conjugates  the \[d$-th power map on the circle to the map \[f\]
on \[J(f)\]. In other words, identifying the unit circle with \[\R/\Z\],
we have a map
$$	\Psi:\R/\Z\;\to\; J(f)	$$
which satisfies \[\Psi(dt)=f(\Psi(t))\], where \[d\] is the degree. Let
\[X\] be the compact set \[\Psi^{-1}(z_0)\subset\R/\Z\]. Since \[z_0\]
is fixed by \[f\], it follows immediately that the map\break \[t\mapsto td
(\mod 1)\] carries \[X\] homeomorphically onto itself. Hence \[X\] is
finite by 18.8, and the conclusion follows.\QED

We will prove the parabolic case of Theorem 18.2 in the following slightly
sharper form.
For the proof in the repelling case, the reader is referred to Petersen.

{\QP{\bf 18.9. Theorem.} \it Let \[f\] be a polynomial of degree \[d\ge 2\]
with  \[K(f)\]  connected,
and let  \[z_0\]  be a parabolic fixed point
with multiplier \[\lambda=1\]. Then for each repelling petal
\[P\]  at  \[z_0\]  there exists at least one
external ray \[R_t\] which lands at \[z_0\] through \[P\], 
where the angle \[t\] is a rational number of the form \[m/(d-1)\],
so that \[R_t\] is mapped onto itself by \[f\].\medskip}

Very roughly, we will show that
the basin \[\C\ssm K\] of the point at infinity corresponds to a union
of one or more annuli in the \'Ecalle cylinder \[P/f\]. Each of these
annuli contains a unique simple closed geodesic \[\Gamma\] (think of
placing a rubber band around a napkin ring); and \[\Gamma\]
will lift to the required external ray in \[\C\ssm K\]. Before giving
details, let us look at two examples.
\smallskip

{\bf Example 1.} Consider the map \[f(z)=z^2+\exp(2\pi i\cdot 3/7)z\] of
Figure 8. Here the seventh iterate is a map of degree \[2^7=128\] with
seven repelling petals about the origin. Hence there must be at least
seven external rays landing at the origin, and their angles must be
rational numbers of the form \[m/127\], so as to map into themselves
under multiplication by 128 modulo 1. In fact a little experimentation
shows that only the ray with angle \[21/127\] and its successive iterates
under doubling modulo 1 will fit in the right order around the origin.
(Compare Goldberg.) Thus there are just seven rays which land at zero,
one in each repelling petal. The numerators of the corresponding angles
are \[21\,,\,37\,,\, 41\,,\, 42\,,\,74\,,\,
82\,,\,84\].\smallskip

{\bf Example 2.} Now consider the cubic map \[f(z)
=z^3-iz^2+z\] of Figure 17. Evidently there is only one repelling petal
at the parabolic fixed point \[z=0\]. Yet we saw above (following 18.4)
that the two rays \[R_0\]
and \[R_{1/2}\] must both land at this point.\smallskip

The proof of 18.9 begins as follows.
According to \S7.7, there is an essentially unique map \[\beta=\alpha^{-1}\]
which is defined and univalent in some left half-plane \[{\cal R}(w) < c\],
taking
values in the repelling petal \[ P\], and which satisfies the Abel equation
\[\beta(w+1) = f(\beta(w))\] or equivalently:
$$	\beta(w) = f(\beta(w-1)) \,. $$
As in 7.11, we can extend to an analytic map from  \[\C\]  to  \[\C\]
satisfying this same equation
by setting  \[\;\beta(w) = f^{\circ n}(\beta(w-n))\;\] for large  \[n\].
{\bf Definition:} Let  \[K_P\]  be the inverse
image of  \[K = K(f)\]  under  \[\beta\]. Since \[K\] is closed and
\[f$-invariant, it follows that \[K_P\]  is closed and periodic:
\[K_P = K_P+1\]. Furthermore, it is not hard to show that \[K_P\] contains
all points \[w=u+iv\] with \[|v|\] large. For points which are
far from the real axis must correspond, under \[\beta\], to points
which belong to one of the two neighboring attracting petals of \[f\].

{\QP{\bf 18.10. Yoccoz Lemma.} \it Each component \[U\] of \[C\ssm K_P\]
is a universal covering of \[C\ssm K\] with projection map \[\beta\].\medskip}
\eject

{\bf Proof.} Consider the topological space \[E\] consisting of all
left infinite sequences \[{\bf w}=(\,\ldots\,,\,w_{-2}\,,\,w_{-1}\,,\,w_0)\] of points in the
Fatou set \[\C\ssm J\] satisfying
$$	\cdots\buildrel f\over\mapsto w_{-2}\buildrel f\over\mapsto w_{-1}
	\buildrel f\over\mapsto w_0	$$
and \[\;\lim_{\,n\to-\infty}\, w_n\,=\,z_0\]. We can give each component of
\[E\] the structure of a Riemann surface in such a way that the projection \[
\pi:(\,\ldots\,,\, w_{-1}\,,\,w_0)\mapsto w_0\] is a local
conformal isomorphism.
In fact, for any such \[{\bf w}\], we can choose
\[N<0\] so that \[w_n\] belongs to some fixed repelling petal \[P\] for \[n
\le N\]. The coordinate \[w_N\in P\] then serves as a local uniformizing
parameter for \[E\], and \[\pi\] corresponds to the map \[f^{|N|}:w_N\mapsto
w_0\]. (It is important here that there are no critical points
in the finite part of the Fatou set.)

To show that \[\pi:E\to \C\ssm K\] is actually a covering map, we
start with any simply connected neighborhood \[V\] of \[\pi({\bf w})=w_0\] in
\[\C\ssm K\]. Then we can find a unique single valued branch \[f_0^{-n}\]
of \[f^{-n}\] on \[V\] which maps \[w_0\] to \[w_{-n}\]. These maps
\[f_0^{-n}\] must constitute a normal family on \[V\], since they take
values in \[\C\ssm K\]. Since
the sequence \[f_0^{-n}\] converges uniformly to \[z_0\] throughout some
small neighborhood of \[w_0\], it follows easily that this sequence
converges to \[z_0\] throughout the given neighborhood \[V\]. Thus for each
\[w\in V\] the sequence \[\;{\bf w}(w)=\big(\,\cdots\,,\, f_0^{-2}(w)\,,\,
f_0^{-1}(w)\,,\,w\big)\;\] belongs to \[E\], hence \[V\] is evenly covered
under the projection \[\pi\].

It is not difficult to check
that \[\C\ssm K_P\] embeds naturally as an open and closed subset of
our covering space \[E\]. But each component \[U\] of \[\C\ssm K_P\] is simply
connected. For otherwise, there would be a compact component of \[\C\ssm U\]
which, after left translation by a large integer, would embed into a compact
of \[K\] in the petal \[P\]. This is impossible, since \[K\] is connected.
Thus each such \[U\] is a universal covering surface of \[\C\ssm K\].\QED

{\bf Remark.} Although \[U\] is a universal covering of \[\C\ssm K\], the proof
does not tell us just how the free cyclic group of deck transformations acts
on \[U\]. We will show in 18.12
that there is also a natural action of the free cyclic
group \[\Z\] by integer translations of \[U\], but that action
is quite different, and has a quotient surface \[U/\Z\] which is an annulus,
naturally embedded in the \'Ecalle cylinder \[P/f\].\smallskip

Let \[\hat G=G\circ\beta\] be the potential function on \[\C\ssm K\] lifted to \[\C\ssm K_P\], so
that\break
\[\hat G(w+1) = \hat G(w)d\] where \[d\] is the degree. Let \[\hat G_0\] be the maximum
value of \[\hat G\] on the imaginary axis, attained say at \[w_0\]. Then
\[\hat G(w)\le G_0\] for \[{\cal R}(w)\le 0\] by the maximum modulus principle
for harmonic maps, hence
$$	 G(w) \le d^n\hat G_0 \qquad{\rm  whenever}\qquad   {\cal R}(w)
 \le n .\eqno	(*)$$

{\QP{\bf 18.11. Lemma.} \it The real part of \[w\] must take values tending
to \[\infty\] within each component \[U\].\medskip}

\noindent This follows, since \[\hat G\] must be unbounded on each \[U\]
by 18.10.\QED

{\QP{\bf 18.12. Main Lemma.} \it Each such component \[U\] is periodic:
\[\;U = U+1\].\medskip}

{\bf Proof.}
Otherwise the translates \[U-n\] would have to be pairwise disjoint. First
consider the component \[U\] on which the maximum \[\hat G_0\] is attained.
Let \[k_n\] be the vertical width of the intersection of \[U-n\] with the
imaginary axis. Then the sum of \[k_n\] is finite, so \[k_n\] tends to
zero as \[n \to +\infty\]. But \[k_n\] is also the vertical width of the
intersection of \[U\] with the line \[\;{\cal R}(w) = n\,\]. {\it Since these
numbers tend to zero, the \Poincare distance of the imaginary axis from
\[\;{\cal R}(w) = n\;\] within \[U\] must grow more than linearly with \[n\].}
\smallskip

On the other hand, an orthogonal trajectory of the curves \[\;\hat G
= {\rm constant}\,\],
of \Poincare arclength \[n \log(d)\] within \[U\], will get from \[w_0\] on the
imaginary axis to a point \[w_n\] with
\[\hat G(w_n) = (d^n)\hat G_0\]. By $(*)$ above, the real
part of \[w_n\] is \[\ge n\]. Thus this distance grows at most linearly
with \[n\]. This contradiction proves the Main Lemma 18.12 when \[w_0\] belongs
to \[U\]. To get a proof for arbitrary \[U\], we must simply use the union of
translates \[U-n\] in place of the entire open set \[C\ssm K_P\] in the
argument above. This proves 18.12.\QED

{\bf Proof of Theorem 18.9.}
We know that each component \[U\] of \[\C\ssm K_P\]
is a universal covering of \[\C\ssm K\], and that the unit translation \[w\mapsto
w+1\] of \[U\] corresponds to the map \[f\] on \[\C\ssm K\].
Furthermore, it follows easily from Corollary 7.5 that the complement
\[\C\ssm K_P\] is non-vacuous, so that there exists at least one such component
\[U\]. Since the imaginary
coordinate is bounded throughout \[U\], the quotient \[U/\Z\] is clearly
an annulus. Hence it has a unique simple closed geodesic with respect to
the Poincar\'e metric. (Compare Problem 2-3.
We can model the annulus as the upper half-plane
\[H\] modulo the identification \[w\sim t_0w\] for some fixed \[t_0>1\]. The
imaginary axis in \[H\] then covers the unique closed geodesic.)\smallskip

This closed geodesic \[\Gamma\subset U/\Z\] lifts to an infinite
\[\Z$-invariant geodesic \[\tilde\Gamma\subset U\], and hence to an
infinite \[f$-invariant geodesic \[\Gamma'=\beta(\tilde\Gamma)\] in
\[\C\ssm K\]. In the negative direction, this geodesic \[\Gamma'\]
converges to the parabolic
fixed point \[z_0\]. For \[\beta\] was constructed in such a way that for
any compact subset \[L\subset U\] the images \[\beta(L-n)\] converge to \[z
_0\] as \[n\to\infty\]. On the other hand, in the positive direction
\[\Gamma'\] converges to the point at infinity, since every compact subset
\[L\subset U\] has the property that the images \[\beta(L+n)=f^{\circ n}
\circ\beta(L)\] converge to \[\infty\] as \[n\to\infty\]. But a \Poincare
geodesic in \[\C\ssm K\] which leads to the point at infinity is necessarily
an external ray.\QED
\vfill

%\rightline{jwm, version of 9-4-91}
\end
